% Gabarit pandoc des corrigés publiés (brut et guidé). Contenu sans marque ni nom de site.
% Compilation : lualatex (latexmk -lualatex).
% Le Markdown reste la source unique : ne pas éditer ce fichier .tex à la main.
\documentclass[11pt,a4paper]{article}

% --- Moteur, polices (Latin Modern) et langue ---------------------------------
\usepackage{iftex}
\ifLuaTeX\else\errmessage{Ce document se compile avec lualatex}\fi
\usepackage{amsmath,amssymb,mathtools}
\usepackage{fontspec}
% mathrm=sym : \mathrm prend ses glyphes dans la police mathématique (ℓ, μ, Ω), comme KaTeX.
\usepackage[math-style=TeX,bold-style=TeX,mathrm=sym]{unicode-math}
% Repli pour les glyphes absents de Latin Modern (indices et exposants Unicode,
% flèches, filets des schémas en texte, symboles) : DejaVu, puis Latin Modern Math
% (ℓ, lettres mathématiques Unicode des légendes).
\directlua{
  luaotfload.add_fallback("repliTexte", {"DejaVuSerif:mode=harf;", "DejaVuSans:mode=harf;", "file:latinmodern-math.otf:mode=harf;"})
  luaotfload.add_fallback("repliMono", {"DejaVuSansMono:mode=harf;", "DejaVuSans:mode=harf;", "file:latinmodern-math.otf:mode=harf;"})
}
\setmainfont{Latin Modern Roman}[RawFeature={fallback=repliTexte}]
\setsansfont{Latin Modern Sans}[RawFeature={fallback=repliTexte}]
\setmonofont{Latin Modern Mono}[RawFeature={fallback=repliMono}]
\setmathfont{Latin Modern Math}
\usepackage[french]{babel}
\frenchsetup{SmallCapsFigTabCaptions=false}

% --- Mise en page --------------------------------------------------------------
\usepackage[a4paper,margin=2.2cm,headheight=14pt,footskip=1.2cm]{geometry}
\usepackage{microtype}
\setlength{\emergencystretch}{4em}
\tolerance=1500
\setlength{\parindent}{0pt}
\setlength{\parskip}{0.55\baselineskip plus 2pt minus 1pt}
\usepackage{xcolor}
% Gris nommés : babel-french rend « ! » actif dans le document.
\definecolor{gris2}{gray}{0.98}
\definecolor{gris3}{gray}{0.97}
\definecolor{gris5}{gray}{0.95}
\definecolor{gris35}{gray}{0.65}
\definecolor{gris40}{gray}{0.60}
\definecolor{gris45}{gray}{0.55}
\definecolor{gris50}{gray}{0.50}
\definecolor{gris55}{gray}{0.45}
\definecolor{gris60}{gray}{0.40}
\definecolor{gris65}{gray}{0.35}
\definecolor{gris70}{gray}{0.30}
\definecolor{gris80}{gray}{0.20}
\definecolor{gris85}{gray}{0.15}
\usepackage{graphicx}
\usepackage{float}
\usepackage{caption}
\captionsetup{font=small,labelformat=empty,justification=centering}
\usepackage{booktabs,longtable,array,multirow,calc}
\usepackage{siunitx}
\sisetup{locale=FR}
\usepackage[version=4]{mhchem}
\usepackage{listings}
\usepackage{fancyhdr}
\usepackage{lastpage}
\usepackage[most]{tcolorbox}
\usepackage{enumitem}
\setlist{nosep,leftmargin=*,topsep=2pt}
\usepackage[hidelinks,unicode]{hyperref}
\usepackage{bookmark}
\hypersetup{pdftitle={CCINP MP-MPI-PC-PSI-TPC-TSI — Français-philosophie 2026},pdfauthor={Rédigé par IA},pdfcreator={LaTeX},pdfsubject={Corrigé guidé}}
\usepackage{xurl}

% --- Pandoc : listes, tableaux, figures ----------------------------------------
\providecommand{\passthrough}[1]{#1}
\providecommand{\tightlist}{\setlength{\itemsep}{0pt}\setlength{\parskip}{0pt}}

\makeatletter
\def\maxwidth{\ifdim\Gin@nat@width>\linewidth\linewidth\else\Gin@nat@width\fi}
\def\maxheight{\ifdim\Gin@nat@height>0.55\textheight 0.55\textheight\else\Gin@nat@height\fi}
\makeatother
\setkeys{Gin}{width=\maxwidth,height=\maxheight,keepaspectratio}
\setcounter{secnumdepth}{-\maxdimen}
\providecommand{\pandocbounded}[1]{#1}
\usepackage{longtable,booktabs,array}

% --- Code (listings) -----------------------------------------------------------
\lstset{
  basicstyle=\ttfamily\small,
  columns=fullflexible,
  keepspaces=true,
  breaklines=true,
  breakatwhitespace=false,
  postbreak=\mbox{\textcolor{gris50}{$\hookrightarrow$}\space},
  frame=single,
  rulecolor=\color{gris35},
  framesep=4pt,
  xleftmargin=5pt,xrightmargin=5pt,
  backgroundcolor=\color{gris3},
  keywordstyle=\bfseries,
  commentstyle=\color{gris65},
  stringstyle=\color{gris80},
  showstringspaces=false,
  tabsize=4,
  extendedchars=true,
  aboveskip=6pt,belowskip=6pt,
}

% --- Macros des corrigés (KaTeX ↔ LaTeX) ---------------------------------------
% Les corrigés sont écrits pour KaTeX : on définit ici ce que KaTeX connaît et
% que LaTeX n'a pas, sans jamais retoucher le Markdown.
\providecommand{\R}{\mathbb{R}}
\providecommand{\N}{\mathbb{N}}
\providecommand{\Z}{\mathbb{Z}}
\providecommand{\Q}{\mathbb{Q}}
\providecommand{\C}{\mathbb{C}}
\providecommand{\dd}{\mathrm{d}}
\providecommand{\llbracket}{\lBrack}
\providecommand{\rrbracket}{\rBrack}
\providecommand{\lBrace}{\lbrace}
\providecommand{\rBrace}{\rbrace}
\providecommand{\Box}{\square}
\providecommand{\Diamond}{\lozenge}
\providecommand{\varnothing}{\emptyset}
\providecommand{\argmax}{\operatorname*{arg\,max}}
\providecommand{\argmin}{\operatorname*{arg\,min}}
\renewcommand{\leq}{\leqslant}
\renewcommand{\geq}{\geqslant}

% Équation hors texte : réduite à la largeur de la ligne si elle déborde.
\newsavebox{\boiteequation}
\newcommand{\equationajustee}[1]{%
  \sbox{\boiteequation}{$\displaystyle #1$}%
  \ifdim\wd\boiteequation>\linewidth
    \[\resizebox{\linewidth}{!}{\usebox{\boiteequation}}\]%
  \else
    \[#1\]%
  \fi}

% --- Encadrés pédagogiques (sobres, lisibles en noir et blanc) -----------------
\tcbset{
  encadre/.style={
    enhanced,breakable,sharp corners,
    colback=white,colframe=gris55,
    boxrule=0.4pt,leftrule=2.2pt,
    left=6pt,right=6pt,top=3pt,bottom=3pt,
    before skip=6pt,after skip=6pt,
    fonttitle=\bfseries\small,coltitle=black,
    attach title to upper,after title={\quad},
    parbox=false,
  }
}
\newtcolorbox{indication}[1]{encadre,colframe=gris45,colback=gris2,title={Indication~#1}}
\newtcolorbox{methode}{encadre,colframe=gris70,title={Méthode}}
\newtcolorbox{piege}{encadre,colframe=gris85,boxrule=0.8pt,leftrule=2.2pt,title={Piège}}
\newtcolorbox{jury}{encadre,colframe=gris60,colback=gris5,title={Jury}}
\newtcolorbox{encadrenote}{encadre,colframe=gris40}

% --- Titres de section ---------------------------------------------------------
\usepackage{titlesec}
\titleformat{\section}{\Large\bfseries}{}{0pt}{}
\titleformat{\subsection}{\large\bfseries}{}{0pt}{}
\titleformat{\subsubsection}{\normalsize\bfseries}{}{0pt}{}
\titlespacing*{\section}{0pt}{14pt plus 4pt minus 2pt}{6pt}
\titlespacing*{\subsection}{0pt}{10pt plus 3pt minus 2pt}{4pt}
\titlespacing*{\subsubsection}{0pt}{8pt plus 2pt minus 2pt}{2pt}

% --- En-tête et pied de page ---------------------------------------------------
\pagestyle{fancy}
\fancyhf{}
\fancyhead[L]{\small\textsc{Corrigé guidé} · CCINP MP-MPI-PC-PSI-TPC-TSI}
\fancyhead[R]{\small Français-philosophie 2026}
\fancyfoot[L]{\footnotesize Corrigé rédigé par IA, non officiel}
\fancyfoot[R]{\footnotesize \thepage\,/\,\pageref*{LastPage}}
\renewcommand{\headrulewidth}{0.4pt}
\renewcommand{\footrulewidth}{0pt}
\fancypagestyle{plain}{\fancyhf{}\fancyfoot[R]{\footnotesize \thepage\,/\,\pageref*{LastPage}}\fancyfoot[L]{\footnotesize Corrigé rédigé par IA, non officiel}\renewcommand{\headrulewidth}{0pt}}

\begin{document}
\thispagestyle{plain}
\begin{center}
  {\small\scshape Corrigé guidé}\par\smallskip
  {\LARGE\bfseries CCINP MP-MPI-PC-PSI-TPC-TSI — Français-philosophie 2026\par}\medskip
  {\large CCINP · MP-MPI-PC-PSI-TPC-TSI · Français-philosophie · 2026\par}
\end{center}
\vspace{2pt}\hrule\vspace{8pt}

\emph{Corrigé rédigé par IA, non officiel. Signalez toute erreur.}

\hypertarget{pruxe9sentation-du-sujet}{%
\section{Présentation du sujet}\label{pruxe9sentation-du-sujet}}

\textbf{L'épreuve.} Épreuve FP4, commune aux filières MP, MPI, PC, PSI, TPC et TSI et mutualisée avec e3a-Polytech : 4 heures, sans document ni machine. Deux exercices évalués indépendamment mais « formant un ensemble cohérent » : un \textbf{résumé} (20 points, un tiers de la note) et une \textbf{dissertation} (40 points, deux tiers). Contraintes matérielles : résumé en \textbf{100 mots ± 10 \%} (donc de 90 à 110 mots) sur la page 2 du Document Réponse, un mot par trait pointillé, changements de paragraphe marqués par une double barre verticale ; début de la dissertation sur les pages 3 et 4 du Document Réponse ; le mot FIN à la fin de la composition. L'énoncé précise que le respect du nombre de mots « représente une part significative du barème » du résumé et que la présentation, l'orthographe et la clarté comptent pour une part importante dans toute la copie. Le rapport 2026 parle, lui, de 10 points pour le résumé et 20 pour la dissertation : la proportion d'un tiers et deux tiers est la même. Il rappelle aussi l'ordre de travail : « L'exercice du résumé doit donc toujours être traité en premier. » (p.~1), car la citation à discuter s'éclaire de l'ensemble du texte.

\textbf{Le texte.} Kathleen Raine (1908-2003), poétesse britannique formée aux sciences naturelles, défend dans \emph{Le monde vivant de l'imagination} une connaissance poétique de la nature. L'extrait (environ 905 mots selon la règle de l'énoncé, quatre paragraphes) suit la marche d'un plaidoyer :

\begin{longtable}[]{@{}
  >{\raggedright\arraybackslash}p{(\columnwidth - 4\tabcolsep) * \real{0.1194}}
  >{\raggedright\arraybackslash}p{(\columnwidth - 4\tabcolsep) * \real{0.0925}}
  >{\raggedright\arraybackslash}p{(\columnwidth - 4\tabcolsep) * \real{0.7581}}@{}}
\toprule\noalign{}
\begin{minipage}[b]{\linewidth}\raggedright
Paragraphe
\end{minipage} & \begin{minipage}[b]{\linewidth}\raggedright
Lignes
\end{minipage} & \begin{minipage}[b]{\linewidth}\raggedright
Mouvement
\end{minipage} \\
\midrule\noalign{}
\endhead
\bottomrule\noalign{}
\endlastfoot
1 & 1-14 & Concession puis thèse : la science quantifie légitimement, mais ne peut se dire \emph{la} connaissance de la nature ; l'art et la simple vision (au sens de contemplation) saisissent des valeurs immensurables ; la nature abrite notre existence spirituelle autant que physique ; conséquence pratique, chacun, contemplateur autant que consommateur, peut s'opposer à la destruction « fonctionnaliste » de la terre. \\
2 & 15-23 & Objection de principe : la science observe le tout depuis la partie ; elle produit une connaissance « objective », mais la prétention de la partie à juger le tout, liée au matérialisme et non inhérente à la méthode, est présomptueuse et rétrécit notre conscience. \\
3 & 24-49 & Appel au retour : revenir au savoir premier, d'avant la séparation de l'intérieur et de l'extérieur et d'avant l'usage ; preuve autobiographique (le groseillier en fleurs, la fleur perçue quintuple avant de savoir compter) ; « il ne s'agissait pas d'apprendre, mais seulement de voir » ; la vie a éloigné l'auteure de cette connaissance. \\
4 & 50-58 & Définition : ce savoir des premières fois n'est pas acquis, c'est une reconnaissance totale et immédiate ; « nous n'avons pas besoin de les connaître pour les reconnaître » ; la nature parle sa propre langue. \\
\end{longtable}

Le rapport 2026 regroupe ces quatre paragraphes en \textbf{trois idées}, les deux derniers formant un seul mouvement (p.~4), et en \textbf{deux temps} : un passage conceptuel et argumentatif (§ 1-2), puis un passage autobiographique, à la première personne, où l'argumentation se fait plus discrète (§ 3-4) (p.~5). Le résumé proposé suit ce découpage en trois paragraphes.

\textbf{Le thème et les œuvres.} Le programme de l'année porte sur l'expérience de la nature : \emph{Vingt mille lieues sous les mers} de Jules Verne (le regard d'un naturaliste, Aronnax, face à un monde inconnu), \emph{La Connaissance de la vie} de Georges Canguilhem (la connaissance du vivant par un vivant, ses méthodes et ses normes) et \emph{Le Mur invisible} de Marlen Haushofer (une femme isolée par une paroi transparente, qui apprend à vivre dans la montagne avec ses bêtes). Le texte de Raine rencontre les trois œuvres sur une même question : la nature se connaît-elle par le savoir ou par une saisie première ?

\textbf{Difficulté.} Le texte est clair mais lyrique et répétitif : le piège du résumé est de suivre ses exemples au lieu de ses idées, et de durcir sa thèse en antiscience. La citation de dissertation, prise dans le texte, est radicale ; elle appelle une discussion réelle, et le mot « reconnaissance » contient un paradoxe exploitable.

\textbf{Ce qui a classé les copies (rapport 2026).} Le jury juge le texte sans difficulté de compréhension, mais relève des points qui ont départagé les candidats :

\begin{itemize}
\tightlist
\item
  au résumé, la \textbf{concession initiale} (« La première idée du texte n'a pas été perçue dans toute sa finesse par la majorité des candidats. », p.~5), le sens de « vision » (contemplation) et surtout celui d'« objective », qui a donné lieu à un contresens « dans un nombre important de copies » (p.~5) ; puis le raisonnement du deuxième paragraphe, la hiérarchie des idées et les calques ;
\item
  à la dissertation, l'\textbf{analyse de la citation} : le pronom « en », le sens d'« appréhender », le concept de « réalité », les adjectifs « totale » et « immédiate », le paradoxe d'une reconnaissance de ce qu'on voit pour la première fois ; puis un plan dont la première partie démontre la pertinence de la comparaison au lieu de traiter un terme par partie.
\end{itemize}

Ce sont les étapes 1 du résumé et 1 à 3 de la dissertation ci-dessous.

\textbf{Gestion du temps.} Suivre le barème : environ 1 h 15 pour le résumé (lecture, découpage, rédaction, deux comptages, recopie), 2 h 30 pour la dissertation (dont 45 minutes de préparation au brouillon), 15 minutes de relecture.

\hypertarget{exercice-1-ruxe9sumuxe9-de-texte-20-points}{%
\section{Exercice 1 --- Résumé de texte (20 points)}\label{exercice-1-ruxe9sumuxe9-de-texte-20-points}}

Résumer, c'est réécrire le texte en réduction, dans son ordre, avec son énonciation, sans commentaire ni citation. Le rapport de réduction est ici d'environ 1 à 9.

\hypertarget{uxe9tape-1-comprendre-le-texte-et-duxe9gager-sa-structure}{%
\subsection{Étape 1 --- Comprendre le texte et dégager sa structure}\label{uxe9tape-1-comprendre-le-texte-et-duxe9gager-sa-structure}}

\begin{indication}{1}

Repérez les mots de liaison qui ouvrent chaque paragraphe et la fonction de chacun (concession, objection, appel, définition). Le texte commence par une négation : que conteste exactement l'auteure, et que ne conteste-t-elle pas ?

\end{indication}

\begin{indication}{2}

Lisez de près la phrase des lignes 20-23 : à qui Raine attribue-t-elle la présomption qu'elle dénonce ? Distinguez la méthode scientifique et l'attitude qui l'accompagne.

\end{indication}

\begin{indication}{3}

Le troisième paragraphe est long parce qu'il est fait d'exemples. Cherchez la phrase qui en donne la conclusion, puis la dernière phrase, qui dit ce que la vie a fait de cette connaissance.

\end{indication}

\hypertarget{solution}{%
\subsubsection{Solution}\label{solution}}

La thèse : la science, qui mesure, ne peut revendiquer le monopole de la connaissance de la nature ; il existe une connaissance première, non quantifiable, une reconnaissance immédiate que la poésie cherche à retrouver. Le texte s'organise en trois idées :

\begin{enumerate}
\def\labelenumi{\arabic{enumi}.}
\tightlist
\item
  \textbf{Concession et thèse} (l. 1-14) : valeur de la science reconnue, mais son droit à se dire \emph{la} connaissance contesté ; la poésie, les arts et la simple contemplation, accessible à tous, atteignent ce qui ne se mesure pas ; la nature est un refuge spirituel autant que physique. Le refus de sa destruction « fonctionnaliste » est une conséquence secondaire.
\item
  \textbf{Critique de la présomption} (l. 15-23) : la science étudie le tout par la partie et sépare le sujet de l'objet, d'où son objectivité ; mais l'homme, qui n'est qu'une partie, prétend juger le tout, comme si sa raison était seule à penser l'univers. Cette présomption, née avec le matérialisme et \emph{non inhérente} à la méthode, appauvrit l'expérience occidentale de la nature.
\item
  \textbf{Appel au savoir premier et définition} (l. 24-58) : retrouver la connaissance d'avant la coupure entre intérieur et extérieur, d'avant l'usage et la convoitise, à laquelle aspire le poète ; l'enfance en témoigne (voir, non apprendre), la vie en a éloigné l'auteure. Ce savoir n'est pas acquis : c'est une reconnaissance totale et immédiate, et la nature se dit elle-même comme une langue.
\end{enumerate}

\begin{piege}

Faire de Raine une ennemie de la science. La première phrase concède sa valeur ; le deuxième paragraphe sépare la méthode de la présomption matérialiste. Un résumé qui écrit « la science est présomptueuse » ou « la science ne connaît rien » commet un contresens. Reformuler la concession par un calque (« je ne nie pas ») ne vaut guère mieux.

\end{piege}

\begin{piege}

Deux contresens de vocabulaire relevés par le jury. « Vision » désigne la contemplation, non l'observation, qui est au contraire une base de la méthode scientifique. « Objective » n'est pas ironique malgré les guillemets : la science est objective parce qu'elle sépare celui qui étudie de ce qu'il étudie. Écrire que la science est « subjective » renverse le paragraphe.

\end{piege}

\begin{jury}

Le rapport 2026 insiste sur la nuance de l'ouverture : « Kathleen Raine ne condamne pas les sciences : elle conteste à la science le droit de se prétendre seule détentrice d'un savoir absolu sur la nature. » (p.~5). Sur le deuxième paragraphe : « la méthode scientifique était dite objective parce qu'elle sépare le sujet qui étudie (le scientifique) de l'objet étudié (la nature) » (p.~6). Il en reconstitue le raisonnement en quatre temps : la science étudie le tout par la partie ; l'homme n'est qu'une partie de la nature ; il ne peut donc se prononcer sur le tout ; en se croyant seul être doué de raison, l'homme occidental n'a de la nature qu'une expérience limitée (p.~6).

\end{jury}

\hypertarget{uxe9tape-2-hiuxe9rarchiser-et-ruxe9partir-les-mots}{%
\subsection{Étape 2 --- Hiérarchiser et répartir les mots}\label{uxe9tape-2-hiuxe9rarchiser-et-ruxe9partir-les-mots}}

\begin{indication}{1}

Partez des proportions du texte (paragraphes d'environ 228, 123, 419 et 135 mots), puis corrigez-les : un paragraphe d'arguments se réduit moins qu'un paragraphe d'exemples.

\end{indication}

\begin{indication}{2}

Pour le souvenir d'enfance, demandez-vous quelle proposition il sert à prouver ; c'est elle seule qui entre dans le résumé, avec au plus un trait concret.

\end{indication}

\hypertarget{solution-1}{%
\subsubsection{Solution}\label{solution-1}}

Budget retenu, pour viser environ 105 mots, en trois paragraphes :

\begin{itemize}
\tightlist
\item
  premier paragraphe, environ 27 mots : concession, arts, poésie et contemplation ouverte à tous, nature refuge spirituel ;
\item
  deuxième, environ 34 mots : le tout étudié par la partie, la séparation du sujet et de l'objet, l'homme simple partie qui prétend juger le tout, l'orgueil matérialiste qui appauvrit l'expérience occidentale ;
\item
  troisième, environ 44 mots : retour au savoir premier, celui du poète, l'enfant qui voit sans apprendre, l'éloignement, la reconnaissance totale et immédiate, la langue de la nature.
\end{itemize}

Le souvenir du groseillier se réduit à sa fonction : l'enfant saisissait les fleurs sans les avoir apprises. Le refus de la destruction de la nature (l. 13), secondaire, est sacrifié.

On garde l'énonciation du texte : Raine argumente au « nous » et témoigne au « je » ; le résumé conserve ce « je » pour le souvenir, qui a valeur de preuve. On proscrit tout métadiscours (« l'auteure montre que »).

\begin{piege}

Grossir ce qui rappelle le cours. Le jury relève des idées « ajoutées, déformées ou survalorisées alors qu'elles étaient secondaires dans le texte, parce qu'elles étaient présentes dans le programme : protection de la nature (l.13) ; utilitarisme (l.13) ; anthropocentrisme (l.20). » (p.~6). À l'inverse, « L'oubli de quelques idées secondaires n'a pas été sanctionné » (p.~6).

\end{piege}

\begin{jury}

Sur la forme : le résumé « doit être structuré en deux ou trois paragraphes (rarement davantage) » (p.~2), et le jury a pénalisé ceux « émiettant le raisonnement en de trop nombreux paragraphes » (p.~2). Sur la voix : « Nous rappelons qu'un soin particulier doit être accordé à l'énonciation, surtout lorsqu'elle fait, comme ici, partie des stratégies argumentatives du texte. » (p.~5)

\end{jury}

\begin{methode}

Les charnières logiques pèsent peu de mots et portent tout le raisonnement : « mais », « pourtant », « donc », les deux-points explicatifs. Un résumé sans elles devient une liste.

\end{methode}

\hypertarget{uxe9tape-3-ruxe9diger-compter-recopier}{%
\subsection{Étape 3 --- Rédiger, compter, recopier}\label{uxe9tape-3-ruxe9diger-compter-recopier}}

\begin{indication}{1}

Reformulez au lieu de recopier : seules quelques notions sans véritable synonyme peuvent être reprises. « Immensurables », pourtant tentant, n'en fait pas partie : cherchez un équivalent.

\end{indication}

\begin{indication}{2}

Comptez selon la règle de l'énoncé, qui n'est pas celle d'un traitement de texte : l'apostrophe et le tiret séparent les mots (« l'art » = 2, « elle-même » = 2), sauf unités insécables (« aujourd'hui » = 1). Comptez deux fois, paragraphe par paragraphe.

\end{indication}

\hypertarget{solution-2}{%
\subsubsection{Solution}\label{solution-2}}

\textbf{Problématique du texte :} la science, qui mesure, peut-elle se dire la connaissance de la nature, quand une reconnaissance première, antérieure à tout savoir acquis, nous la livre tout entière ?

\begin{quote}

La science mesure utilement la nature, sans en détenir seule la connaissance : arts, poésie et simple contemplation, ouverte à tous, saisissent l'incalculable de ce refuge spirituel. ‖ Étudiant le tout par la partie, séparant observateur et objet, la science est objective. Mais l'homme, simple fragment, prétend juger l'ensemble en seul être pensant : cet orgueil matérialiste appauvrit notre expérience occidentale. ‖ Retrouvons donc, avec le poète, le savoir premier, antérieur à la coupure entre moi et monde. Enfant, je saisissais fleurs et bêtes sans apprendre ; grandir m'en a éloignée. On reconnaît ces êtres d'emblée, entièrement, sans les connaître : la nature parle sa langue.

\end{quote}

\textbf{105 mots} (27 + 34 + 44), dans la fourchette de 90 à 110 ; doubles barres non comptées.

\begin{piege}

Viser la borne (110 ou 90) : la moindre divergence de comptage sur une élision fait sortir de la tolérance. Visez l'intérieur de la fourchette.

\end{piege}

\begin{piege}

Les calques. Le rapport 2026 cite les reprises des termes « immensurable », « valeur », « droit de la science » ou « évaluer la nature » (p.~6) : ce sont les mots que l'on recopie par prudence, et qu'il fallait justement reformuler.

\end{piege}

\begin{jury}

Les quatre critères du rapport 2026 sont le repérage des idées essentielles, une construction « cohérente, fluide et fidèle à la thèse générale du texte », des reformulations personnelles et le respect du nombre de mots (p.~2). On attendait que le résumé « évite les calques (a fortiori les citations ou recopiages !) » (p.~2). Côté comptage, « La quasi-totalité des résumés compte entre 90 et 110 mots. » (p.~5), mais « Quelques candidats commettent encore l'erreur de compter par exemple « l'enfant » comme un seul mot » (p.~5). Enfin, « Très peu de copies ont pris soin d'indiquer, par des doubles barres obliques, le passage d'un paragraphe à un autre. » (p.~5)

\end{jury}

\hypertarget{exercice-2-dissertation-40-points}{%
\section{Exercice 2 --- Dissertation (40 points)}\label{exercice-2-dissertation-40-points}}

Le sujet demande « dans quelle mesure » la lecture des œuvres permet de souscrire au jugement de Raine : il faut donc évaluer ce jugement, ni l'adopter tel quel ni le rejeter en bloc, à partir des trois œuvres.

\hypertarget{uxe9tape-1-analyser-la-citation}{%
\subsection{Étape 1 --- Analyser la citation}\label{uxe9tape-1-analyser-la-citation}}

\begin{indication}{1}

Replacez la phrase dans son contexte (l. 50-54) : à quoi renvoie le pronom « en » ? De quelle réalité parle-t-on ?

\end{indication}

\begin{indication}{2}

Pesez chacun des termes : « expérience acquise par le savoir » (quelle expérience ?), « appréhender », « totale », « immédiate ». Puis regardez le préfixe de « reconnaissance ».

\end{indication}

\hypertarget{solution-3}{%
\subsubsection{Solution}\label{solution-3}}

\begin{itemize}
\tightlist
\item
  Le pronom « en » renvoie à l'énumération des lignes 50-52 : végétaux (arbre en fleur, corolle d'une pensée), animaux (oiseau, faisan, chat), astre, cailloux, ruisseau, et un intrus, des tasses accrochées dans un dressoir. Il s'agit de la réalité des choses, et d'abord de la nature, rencontrées pour la première fois.
\item
  « Appréhender » : saisir par l'esprit, connaître, et non redouter ou affronter.
\item
  « Réalité » : ce que les choses sont en elles-mêmes, par-delà nos représentations ; à ne pas confondre avec la vérité, qui qualifie un jugement.
\item
  L'« expérience acquise par le savoir » n'est pas l'expérience vécue, mais l'expérience instruite : transmise par l'école, les livres, la parole d'autrui, la science aussi, elle nomme, classe, mesure, et s'accumule. Elle est médiate, progressive, et étudie la nature par ses parties.
\item
  « Totale » : portant sur l'être entier, non sur un caractère isolé ; « immédiate » : sans intermédiaire (concept, nom, instrument) et sans délai, de l'ordre de la révélation et de la fusion entre celui qui contemple et ce qu'il contemple.
\item
  « Reconnaissance » : connaître \emph{de nouveau}. Le mot suppose une familiarité antérieure, que Raine appelle un savoir « inné ». D'où un paradoxe : la formule nie que le savoir donne accès au réel tout en supposant un savoir premier.
\end{itemize}

\begin{methode}

Une citation de dissertation se lit dans son contexte quand le texte est fourni : le contexte fixe le sens des pronoms et des termes, et évite de disserter sur un sujet voisin.

\end{methode}

\begin{piege}

Deux contresens lexicaux. « Appréhender » pris au sens de « redouter » conduit à disserter sur l'hostilité de la nature et la survie : selon le rapport 2026, « Cette méconnaissance du vocabulaire a ainsi provoqué les très rares hors-sujets de la session. » (p.~8). « Reconnaissance » prise au sens de « gratitude » : le jeu sur la polysémie est possible dans un dépassement, mais « traiter la reconnaissance comme gratitude dès la première partie relevait du contresens » (p.~8).

\end{piege}

\begin{jury}

Le rapport 2026 fait de l'analyse du sujet le point faible de la session : « Presque aucun candidat n'a analysé ce pronom « en » » (p.~7) ; « « Appréhender » signifiait ici « connaître » ou « saisir par l'esprit ». » (p.~8) ; « le concept de « réalité » n'a pratiquement jamais été ni relevé, ni étudié sérieusement » (p.~8) ; « les adjectifs « totale » et « immédiate » n'ont souvent été ni relevés, ni travaillés » (p.~8). Sur le paradoxe de la reconnaissance d'un enfant qui voit pour la première fois : « Très peu de copies ont relevé ce paradoxe ; aucune n'a été sanctionnée pour cet oubli. » (p.~8). Enfin, l'expérience acquise par le savoir « ne devait pas être réduite à la science, encore moins à l'expérimentation scientifique » (p.~8).

\end{jury}

\hypertarget{uxe9tape-2-probluxe9matiser}{%
\subsection{Étape 2 --- Problématiser}\label{uxe9tape-2-probluxe9matiser}}

\begin{indication}{1}

Opposez ce que la thèse affirme à ce que l'expérience commune et les œuvres suggèrent : la réalité de la nature se donne-t-elle d'emblée, ou s'apprend-elle ?

\end{indication}

\begin{indication}{2}

Exploitez le paradoxe de la \emph{re}-connaissance : si reconnaître suppose d'avoir connu, d'où vient ce premier savoir ? Ne pourrait-il pas être, aussi, le fruit d'un apprentissage ?

\end{indication}

\hypertarget{solution-4}{%
\subsubsection{Solution}\label{solution-4}}

\textbf{Problématique :} la réalité de la nature se donne-t-elle d'emblée, dans une reconnaissance que le savoir acquis ne ferait que recouvrir, ou ne se livre-t-elle qu'au terme d'un apprentissage qui seul permet de la reconnaître ?

\begin{piege}

Réduire le sujet à « la science contre la poésie » ou « la raison contre le cœur ». Raine concède la valeur de la science ; et Canguilhem, au programme, est un philosophe de la \emph{connaissance} de la vie. Le cœur du sujet est l'opposition entre savoir acquis et saisie immédiate. Le rapport 2026 est net : « Trop de copies se sont livrées à une critique radicale de la science, ce qui est préoccupant chez de futurs ingénieurs. » (p.~9)

\end{piege}

\begin{jury}

Le rapport 2026 refuse la simple reformulation : « Trop de candidats se contentent de reformuler la citation sous une forme interrogative » (p.~9), et multiplier les questions ne remplace pas un problème. Il suggère deux problématiques : « Quelles sont les expériences qui donnent accès à la réalité de la nature ? Le savoir fait-il écran à la réalité de la nature ? » (p.~9). Il rappelle aussi la grammaire : au style indirect, pas d'inversion du sujet et un point final ; au style direct, un point d'interrogation (p.~9-10).

\end{jury}

\hypertarget{uxe9tape-3-construire-le-plan-et-choisir-les-ruxe9fuxe9rences}{%
\subsection{Étape 3 --- Construire le plan et choisir les références}\label{uxe9tape-3-construire-le-plan-et-choisir-les-ruxe9fuxe9rences}}

\begin{indication}{1}

Cherchez d'abord, dans chaque œuvre, une situation où le savoir fait écran, puis une situation où l'évidence immédiate se révèle trompeuse, enfin une situation où savoir et regard se rejoignent.

\end{indication}

\begin{indication}{2}

Pour Verne, pensez au trio Conseil, Ned Land, Aronnax et à l'ouverture du roman ; pour Canguilhem, au milieu vécu, au monstre et à l'expérimentation ; pour Haushofer, à ce que la narratrice doit apprendre et à ce qu'elle finit par voir.

\end{indication}

\begin{indication}{3}

La troisième partie ne doit pas être un « il faut nuancer » : elle doit dire précisément ce que la formule a de juste et ce qu'elle outrepasse.

\end{indication}

\hypertarget{solution-5}{%
\subsubsection{Solution}\label{solution-5}}

\begin{longtable}[]{@{}
  >{\raggedright\arraybackslash}p{(\columnwidth - 8\tabcolsep) * \real{0.1540}}
  >{\raggedright\arraybackslash}p{(\columnwidth - 8\tabcolsep) * \real{0.1778}}
  >{\raggedright\arraybackslash}p{(\columnwidth - 8\tabcolsep) * \real{0.2338}}
  >{\raggedright\arraybackslash}p{(\columnwidth - 8\tabcolsep) * \real{0.2209}}
  >{\raggedright\arraybackslash}p{(\columnwidth - 8\tabcolsep) * \real{0.1835}}@{}}
\toprule\noalign{}
\begin{minipage}[b]{\linewidth}\raggedright
Partie
\end{minipage} & \begin{minipage}[b]{\linewidth}\raggedright
Thèse
\end{minipage} & \begin{minipage}[b]{\linewidth}\raggedright
Verne
\end{minipage} & \begin{minipage}[b]{\linewidth}\raggedright
Canguilhem
\end{minipage} & \begin{minipage}[b]{\linewidth}\raggedright
Haushofer
\end{minipage} \\
\midrule\noalign{}
\endhead
\bottomrule\noalign{}
\endlastfoot
I. Le savoir acquis peut faire écran & La formule nomme une expérience réelle : un savoir qui a tous les mots et ne rencontre rien. & Conseil classe sans voir ; Aronnax ébloui devant la vitre du salon, la forêt de l'île Crespo, l'Atlantide. & Introduction : la pensée du vivant procède du vivant ; « Le vivant et son milieu » : le milieu vécu précède le milieu abstrait de la science. & Le savoir d'avant le mur perd son objet ; présence de Lynx, de la forêt ; le mur reste inexpliqué. \\
II. L'immédiat est le lieu de l'illusion & La réalité se conquiert. & Le « monstre » de 1866, évidence unanime et fausse ; l'hypothèse du narval ; la coque de tôle. & « La monstruosité et le monstrueux » : la tératologie dissipe l'imaginaire ; « L'expérimentation en biologie animale » : le détour nécessaire. & L'apprentissage de la survie : cultures, foin, traite, chasse, réserves. \\
III. La reconnaissance comme fruit d'un savoir qui retourne au regard & La formule vaut comme ordre de fondation, non comme exclusion du savoir. & Ned et Conseil complémentaires ; Aronnax voit plus parce qu'il sait ; le massacre des cachalots, savoir devenu pouvoir ; la question de l'Ecclésiaste à la fin. & La vie forme, la connaissance analyse ce que la vie a formé ; « Le normal et le pathologique » : retour au vivant. & Attention instruite ; la corneille blanche ; l'homme à la hache, violence sans regard. \\
\end{longtable}

Chaque œuvre traverse les trois parties : c'est ce qui fait des œuvres des arguments et non des illustrations. Dans le devoir, chaque paragraphe s'ouvre sur l'argument, jamais sur le nom d'un auteur, et des passerelles relient les œuvres deux à deux (le milieu vécu de Canguilhem et la montagne de la narratrice, le « monstre » de 1866 et le monstre biologique).

\begin{piege}

Consacrer la première partie à la reconnaissance immédiate et la deuxième au savoir acquis, une notion par partie. Le jury note que « De nombreuses copies » l'ont fait (p.~10) : on juxtapose alors deux thèmes au lieu de discuter une thèse.

\end{piege}

\begin{jury}

Pour le plan : « Or la thèse aurait dû démontrer la pertinence de cette comparaison, et l'antithèse la discuter. » (p.~10). Pour les paragraphes : « On ne saurait accepter qu'on commence un paragraphe, et moins encore une grande partie, par une référence directe ou un emprunt sec à un auteur ou à une œuvre. » (p.~2). Pour les œuvres, le rapport recommande d'exploiter « des couples ou paires d'œuvres dans chaque argument » (p.~3) plutôt que des allusions systématiques aux trois. Parmi les dépassements qu'il propose figurent l'idée, tirée de Canguilhem, que la connaissance et la vie ne s'opposent pas, et celle d'un savoir qui permet de s'émerveiller « avec l'œil averti du savant » (p.~11) : c'est la troisième partie retenue ici.

\end{jury}

\begin{piege}

Inventer ou approximer une citation. Une scène précisément située (les panneaux du salon, la corneille blanche) vaut mieux qu'une phrase entre guillemets dont on n'est pas sûr ; en cas de doute, on paraphrase sans guillemets.

\end{piege}

\hypertarget{uxe9tape-4-ruxe9diger-lintroduction}{%
\subsection{Étape 4 --- Rédiger l'introduction}\label{uxe9tape-4-ruxe9diger-lintroduction}}

\begin{indication}{1}

Une introduction complète comporte : une amorce, la citation exacte avec son auteur et son œuvre, l'analyse des termes, la problématique, l'annonce du plan avec les trois œuvres.

\end{indication}

\begin{indication}{2}

Cherchez une amorce qui pose déjà le problème : par exemple l'idée traditionnelle que l'étonnement est le commencement du savoir, que Raine renverse.

\end{indication}

\hypertarget{solution-6}{%
\subsubsection{Solution}\label{solution-6}}

L'introduction retenue part de l'idée que l'étonnement est le commencement de la connaissance, que Raine retourne en faisant de l'émerveillement une connaissance à part entière ; elle cite la phrase, en précise le référent (l'énumération des lignes 50-52, intrus compris), fixe le sens d'« appréhender » et de « réalité », définit ses deux termes opposés, dégage le paradoxe de la \emph{re}-connaissance, pose la problématique et annonce les trois parties en nommant les trois œuvres. Elle est rédigée en entier dans le devoir ci-dessous.

\begin{methode}

La problématique se formule comme une vraie alternative, dont les deux termes sont défendables. L'annonce de plan doit laisser deviner la réponse finale.

\end{methode}

\begin{jury}

Le rapport 2026 constate que les critères formels de l'introduction sont presque toujours respectés, mais que « Le jury a lu trop d'amorces très longues, inutiles, sans lien avec le sujet » (p.~9), alors que les analyses de la citation restent souvent trop courtes. L'équilibre est donc à inverser : amorce brève, analyse développée.

\end{jury}

\hypertarget{uxe9tape-5-le-devoir-complet}{%
\subsection{Étape 5 --- Le devoir complet}\label{uxe9tape-5-le-devoir-complet}}

\begin{indication}{1}

Chaque paragraphe argumentatif suit le schéma : idée directrice, référence précise et analysée, retour au sujet. Les transitions entre parties résument et relancent.

\end{indication}

\begin{indication}{2}

La conclusion répond à la question posée (« dans quelle mesure ») en une phrase assertive, puis peut ouvrir en reprenant une image du texte.

\end{indication}

\hypertarget{solution-7}{%
\subsubsection{Solution}\label{solution-7}}

\hypertarget{introduction}{%
\paragraph{Introduction}\label{introduction}}

Une longue tradition philosophique place l'étonnement au commencement du savoir : on se met à chercher parce qu'on s'émerveille. Kathleen Raine renverse cette continuité. Dans \emph{Le monde vivant de l'imagination}, la poétesse britannique soutient que l'émerveillement n'est pas l'antichambre de la connaissance, mais une connaissance à part entière, et la plus haute : « Ce n'est pas l'expérience acquise par le savoir qui nous permet d'en appréhender la réalité, mais une reconnaissance totale et immédiate. » Le pronom « en » renvoie à l'énumération qui précède (l. 50-52) : un arbre en fleur, un oiseau, la lune, des cailloux, un ruisseau, un faisan, un chat, la corolle d'une pensée, et jusqu'à un intrus, des tasses accrochées dans un dressoir. Il s'agit donc des choses, et d'abord des êtres de la nature, telles qu'on les rencontre pour la première fois. « Appréhender » ne veut pas dire ici redouter, mais saisir par l'esprit, connaître. Quant à la « réalité », ce n'est pas la vérité d'un énoncé, mais ce que les choses sont en elles-mêmes, par-delà nos représentations. L'« expérience acquise par le savoir » désigne l'expérience instruite : transmise par l'école, les livres et la parole d'autrui autant que par la science, elle nomme, classe et mesure, et s'accumule avec le temps. La « reconnaissance totale et immédiate », à l'inverse, saisit l'être entier, sans médiation ni délai, dans une fusion du regard et de la chose. Le jugement est donc double : il refuse au savoir le pouvoir de nous livrer le réel et l'attribue à une évidence première. Il recèle pourtant un paradoxe que le mot même de Raine trahit : \emph{re}-connaître, c'est connaître de nouveau, et l'immédiateté revendiquée suppose une familiarité antérieure, que l'auteure nomme un savoir « inné ». La réalité de la nature se donne-t-elle d'emblée, dans une reconnaissance que le savoir acquis ne ferait que recouvrir, ou ne se livre-t-elle qu'au terme d'un apprentissage qui seul permet de la reconnaître ? Les trois œuvres du programme, \emph{Vingt mille lieues sous les mers} de Jules Verne, \emph{La Connaissance de la vie} de Georges Canguilhem et \emph{Le Mur invisible} de Marlen Haushofer, accordent d'abord à Raine que le savoir accumulé peut faire écran à la présence des êtres. Elles montrent ensuite que l'évidence immédiate est le lieu même de l'illusion et que la réalité se conquiert. Elles invitent enfin à comprendre la reconnaissance non comme un point de départ perdu, mais comme le fruit d'un savoir qui retourne au regard.

\hypertarget{i.-le-savoir-acquis-peut-faire-uxe9cran-uxe0-la-ruxe9alituxe9-des-uxeatres}{%
\paragraph{I. Le savoir acquis peut faire écran à la réalité des êtres}\label{i.-le-savoir-acquis-peut-faire-uxe9cran-uxe0-la-ruxe9alituxe9-des-uxeatres}}

Le premier mérite de la formule est de nommer une expérience que les trois œuvres attestent : celle d'un savoir qui possède tous les mots et ne rencontre rien.

Le savoir des classifications, d'abord, peut se suffire à lui-même au point de ne plus rien voir, et Verne en fait un type romanesque. Dans \emph{Vingt mille lieues sous les mers}, Conseil, le domestique du professeur Aronnax, connaît par cœur la classification des naturalistes, embranchements, classes, ordres, familles, genres, espèces ; mais le narrateur précise d'emblée que cette science est de pure théorie et qu'il ne distinguerait pas un cachalot d'une baleine. Conseil incarne exactement l'« expérience acquise par le savoir » réduite à elle-même : un classement qui tourne à vide. Devant la vitre du salon du \emph{Nautilus}, lorsque les panneaux s'ouvrent sur la mer éclairée, c'est Ned Land, le harponneur, qui reconnaît et nomme les poissons, tandis que Conseil se contente de les ranger. Et Aronnax lui-même, professeur au Muséum, cesse alors d'être un savant qui énumère : il est un spectateur ébloui, qui avoue à plusieurs reprises, devant la forêt sous-marine de l'île Crespo ou les ruines englouties de l'Atlantide, que ses connaissances ne suffisent pas à dire ce qu'il voit. Le roman fait ainsi de l'émerveillement le moment où le réel déborde le savoir.

Il arrive même que la réalité de la nature ne se découvre qu'au moment où le savoir acquis cesse de servir. \emph{Le Mur invisible} fait de cette épreuve son point de départ. La narratrice, une citadine venue passer quelques jours dans un pavillon de chasse des montagnes autrichiennes, se découvre un matin séparée du monde par une paroi transparente. Tout ce qu'elle savait, les usages sociaux, les nouvelles, les compétences d'une vie urbaine, perd d'un coup son objet. Ce qui s'impose alors ne procède d'aucun savoir : la présence de Lynx, le chien, dont l'attachement se donne sans avoir à être appris ; la forêt, la prairie, l'alpage, qui deviennent la totalité de son monde et qu'elle apprend à regarder pour eux-mêmes. Le mur lui-même reste une énigme qu'elle renonce à expliquer : elle écrit son récit non pour comprendre, mais pour tenir. Haushofer donne ici raison à Raine : la perte des savoirs ordinaires rend la narratrice à la présence des êtres.

Cette antériorité du rapport vécu sur le savoir n'est pas seulement une donnée littéraire : elle a un fondement philosophique. Dans « La pensée et le vivant », introduction de \emph{La Connaissance de la vie}, Canguilhem rappelle que celui qui connaît la vie est lui-même un vivant, et que la pensée du vivant doit tenir du vivant l'idée du vivant : la biologie ne peut se donner son objet de l'extérieur. Le chapitre « Le vivant et son milieu » précise ce primat : le milieu homogène et décentré de la physique est une construction obtenue par abstraction, à partir du milieu vécu qu'un organisme centre autour de lui et qu'il qualifie selon ses besoins. C'est un tel milieu centré, qualifié par les besoins et les attachements, qui devient le monde entier de la narratrice du \emph{Mur invisible}. Même le chapitre consacré à la théorie cellulaire montre que les concepts de la biologie sont nés d'images et de valeurs, avant toute mesure. En ce sens, la reconnaissance dont parle Raine n'est pas un ornement poétique ajouté à la science : elle est ce sans quoi il n'y aurait même pas d'objet à connaître.

Il faut pourtant aller plus loin : si l'évidence première suffisait, aucun de ces livres ne raconterait une histoire, car il ne s'y passerait rien.

\hypertarget{ii.-mais-limmuxe9diat-est-le-lieu-propre-de-lillusion-la-ruxe9alituxe9-se-conquiert}{%
\paragraph{II. Mais l'immédiat est le lieu propre de l'illusion : la réalité se conquiert}\label{ii.-mais-limmuxe9diat-est-le-lieu-propre-de-lillusion-la-ruxe9alituxe9-se-conquiert}}

Loin de livrer la réalité, la perception immédiate est d'abord le règne de la projection. \emph{Vingt mille lieues sous les mers} s'ouvre précisément sur une perception immédiate, totale et fausse. En 1866, des navires de plusieurs nations aperçoivent en pleine mer une chose énorme, longue, fusiforme, parfois phosphorescente : l'opinion y voit aussitôt un monstre marin. Aronnax lui-même, consulté par la presse, propose l'hypothèse d'un narval gigantesque, erreur savante construite sur une évidence populaire. Il faut des mois de traque à bord de l'\emph{Abraham Lincoln} puis la chute du professeur à la mer, qui le fait échouer sur le dos même de la bête, pour que la vérité apparaisse : sous ses pieds, Aronnax découvre une coque de tôle, celle d'un bateau sous-marin. Verne place ainsi, en ouverture, la réfutation de l'immédiateté : la première fois que l'on voit une chose, on la reconnaît d'après ce que l'on craint ou ce que l'on attend, et l'on se trompe. La reconnaissance n'est jamais vierge ; elle mobilise des images reçues.

Ce que la réaction spontanée prend pour le réel, seul un savoir patient peut le rectifier. Canguilhem le montre dans « La monstruosité et le monstrueux » : comme le « monstre » de 1866, le monstre biologique suscite d'emblée une réaction totale, faite de fascination et d'effroi, parce qu'il semble contredire l'ordre même de la vie. Cette réaction est une évaluation spontanée du vivant par le vivant ; elle n'est pas une connaissance. C'est le travail de la tératologie, avec les Geoffroy Saint-Hilaire, que de soustraire la monstruosité au monstrueux en montrant qu'elle relève des mêmes lois du développement que la forme normale. Le savoir acquis, ici, ne recouvre pas le réel : il le dégage de l'imaginaire qui l'enveloppait. De même, le chapitre « L'expérimentation en biologie animale » établit qu'on ne connaît le vivant qu'au prix d'un détour, isoler, provoquer, faire varier, détour d'autant plus délicat que l'organisme est un individu et que l'intervention le modifie, mais dont aucune contemplation ne dispense.

Habiter la nature, enfin, ne relève pas d'un don mais d'un apprentissage, et \emph{Le Mur invisible} ne raconte rien d'autre. La narratrice ne survit pas parce qu'elle reconnaît la nature, mais parce qu'elle apprend, lentement et maladroitement, à planter des pommes de terre et des haricots, à faucher et rentrer le foin, à traire et soigner la vache Bella, à chasser pour nourrir le chien et se nourrir elle-même, à prévoir les réserves de l'hiver. Chaque erreur se paie en fatigue, en récoltes compromises, en bêtes perdues. La montagne n'est pas pour elle un livre ouvert, mais un milieu opaque dont il faut patiemment déchiffrer les saisons, les orages, les dangers. Rien d'inné ne la sauve : c'est une compétence construite, jour après jour, qui lui permet d'habiter ce qu'elle se contentait autrefois de traverser. Le roman répond ainsi à Raine par sa forme même : la réalité de la nature s'y livre dans la durée, saison après saison, et non dans l'instant d'une révélation.

La formule de Raine semble donc démentie : loin d'être un écran, le savoir acquis est souvent la condition pour voir ce qui est. Faut-il pour autant renvoyer l'intuition à l'illusion ?

\hypertarget{iii.-la-reconnaissance-comme-fruit-dun-savoir-qui-retourne-au-regard}{%
\paragraph{III. La reconnaissance comme fruit d'un savoir qui retourne au regard}\label{iii.-la-reconnaissance-comme-fruit-dun-savoir-qui-retourne-au-regard}}

Ce serait manquer ce que la formule a de juste. Elle ne décrit pas une alternative, mais un ordre : le savoir naît d'une rencontre et doit y revenir. Le texte lui-même le laisse voir. Raine a appris plus tard que le buisson de son enfance était un groseillier, et ce nom n'a rien détruit : elle revoit les pétales avec la même précision, et c'est le nom qui lui permet aujourd'hui de nous transmettre sa vision. L'enfant perçoit la fleur comme quintuple avant de savoir compter jusqu'à cinq ; mais le nombre viendra, et il dira exactement ce que le regard avait saisi. La poésie elle-même, que Raine donne pour modèle, n'est pas une saisie immédiate : elle est un travail de la langue, une étude qui se déploie dans la durée. La reconnaissance est première dans l'ordre du fondement ; elle n'est pas suffisante dans l'ordre de la connaissance.

Le savoir n'éteint pas le regard : il l'aiguise, tant qu'il ne se fait pas pouvoir. Verne le montre en distribuant les rôles. Conseil classe sans voir, Ned Land voit sans classer, et le roman suggère qu'à eux deux ils feraient un naturaliste accompli : la reconnaissance véritable naît de leur union. Ned, d'ailleurs, ne voit dans la mer que des proies à harponner : son regard immédiat est celui du « consommateur » dont parle Raine. Aronnax, lui, voit davantage parce qu'il sait : la classification ne s'interpose pas entre lui et l'animal, elle multiplie ce qu'il perçoit, distingue là où l'ignorant ne verrait qu'une masse confuse, et fait de la vitre du salon un spectacle inépuisable. Le roman indique aussi le point où ce savoir se pervertit : lorsque le \emph{Nautilus}, armé de son éperon, massacre une troupe de cachalots, le récit insiste sur le carnage et la mer ensanglantée. Le savoir aveugle n'est pas celui qui sait trop, mais celui qui devient pouvoir. Et c'est l'expérience acquise, au terme d'un voyage de vingt mille lieues, qui permet à Aronnax de conclure qu'à la vieille question de l'Ecclésiaste, qui a jamais sondé les profondeurs de l'abîme, Nemo et lui sont désormais les seuls à pouvoir répondre.

Connaître et vivre ne s'opposent donc pas : la connaissance prolonge la vie, à condition de lui revenir. Canguilhem donne à cet ordre sa formule dans l'introduction de son livre : la vie forme des formes, la connaissance analyse les matières que la vie a informées. La connaissance n'est pas le contraire de la vie ; elle en est un moment, qui défait pour comprendre ce que la vie a formé, et qui doit donc toujours revenir à ce qu'elle a défait. Le chapitre « Le normal et le pathologique » en tire une conséquence concrète : la norme n'est pas une moyenne statistique, elle est instituée par le vivant lui-même, capable de créer de nouvelles normes ; un savoir médical qui oublierait l'expérience vécue du malade manquerait son objet. Il n'y a donc pas à choisir entre reconnaître et connaître, mais à exiger du savoir qu'il n'oublie pas d'où il vient.

La reconnaissance la plus haute, enfin, vient à la fin du chemin et non à son début. C'est ce que la narratrice de Haushofer atteint au terme de son épreuve. Après des années de travail, son regard n'est plus celui de la femme d'avant, ni celui d'une enfant qui reconnaîtrait tout d'emblée : c'est une attention instruite, faite d'habitude, de soin et de deuils. Elle finit par se sentir partie de la forêt plutôt que sujet face à elle, comme si la séparation du dedans et du dehors dont parle Raine s'effaçait, non par un retour à l'enfance, mais au bout d'un long travail. Elle remarque une corneille blanche que ses congénères rejettent parce qu'elle ne leur ressemble pas, et elle la nourrit : reconnaître un être singulier dans sa fragilité n'est pas ici un don originel, mais le fruit tardif d'une longue fréquentation. À l'inverse, l'homme qui surgit à la fin du récit et abat à coups de hache le jeune taureau puis Lynx incarne une violence sans regard ni attention. Chez Haushofer, l'immédiat sans savoir n'est pas le paradis perdu : c'est la brutalité.

\hypertarget{conclusion}{%
\paragraph{Conclusion}\label{conclusion}}

On peut donc souscrire au jugement de Kathleen Raine, mais dans une mesure précise. Il est vrai qu'aucune accumulation de savoirs ne produit à elle seule la présence des êtres et que le savoir peut tourner à vide : Conseil le montre, et Canguilhem le fonde en rappelant que la pensée du vivant procède du vivant. Il est faux, en revanche, que la reconnaissance immédiate suffise à saisir la réalité : le monstre de 1866, le monstrueux qu'a dû dissiper la tératologie, la montagne opaque de la narratrice sont autant d'évidences trompeuses ou de mondes à déchiffrer. La formule vaut comme ordre de fondation et comme avertissement contre la réduction de la nature à une matière exploitable ; elle ne vaut pas comme exclusion du savoir. Les trois œuvres décrivent un même trajet : on ne voit d'abord que ce que l'on projette, puis on apprend, et le savoir, s'il reste fidèle à la vie, rend enfin au regard ce qu'il lui avait d'abord pris. Si, comme l'écrit Raine, la nature parle une langue, il faut lui accorder ce que l'on accorde à toute langue : on la reconnaît d'autant mieux qu'on l'a apprise.

FIN

\begin{piege}

L'exemple raconté. Conseil qui classe, Aronnax ébloui, l'attachement à Lynx sont des exemples « attendus voire convenus » selon le rapport 2026 (p.~11) : ils ne valent que s'ils sont interprétés. Autre faute relevée : les noms d'auteurs mal orthographiés (Verne, Canguilhem, Haushofer), jugés inadmissibles après une année d'étude (p.~7).

\end{piege}

\begin{jury}

Le rapport 2026 sur les exemples : « Un bon exemple n'est jamais narratif ni descriptif ; il est précis et interprété. » (p.~11). Sur les transitions : « Les transitions peuvent être courtes, mais elles sont nécessaires à la fin de chaque grande partie. » (p.~11). Sur la conclusion, qui doit s'engager pour une thèse : « On pourra accepter la conclusion dite ouverte si et seulement si elle ne se termine pas par une interrogation passe-partout. » (p.~3)

\end{jury}

\hypertarget{bilan}{%
\section{Bilan}\label{bilan}}

\begin{itemize}
\tightlist
\item
  \textbf{Résumé.} Suivre les idées, pas les exemples ; garder les charnières logiques ; respecter les nuances de l'auteur, en particulier les concessions (ici, la valeur de la science et la distinction entre la méthode et la présomption matérialiste) ; compter selon la règle de l'énoncé, deux fois, et viser l'intérieur de la fourchette.
\item
  \textbf{Dissertation.} Relire la citation dans son contexte ; peser chaque terme et chercher le mot qui porte un paradoxe (ici, « reconnaissance ») ; formuler une problématique qui soit une vraie alternative ; construire une troisième partie qui dit ce que la thèse a de juste et ce qu'elle outrepasse.
\item
  \textbf{Œuvres.} Préparer pour chaque œuvre des scènes précisément situées, réutilisables sur tout sujet du thème : pour Verne, Conseil et Ned Land, la vitre du salon, l'ouverture sur le « monstre », les cachalots, la fin du roman ; pour Canguilhem, l'introduction, le milieu, l'expérimentation, le normal et le pathologique, le monstre ; pour Haushofer, le mur, les travaux et les saisons, les bêtes, la corneille blanche, l'homme à la hache.
\item
  \textbf{Questions discriminantes en 2026.} D'après le rapport, la concession initiale et le sens d'« objective » au résumé ; l'analyse du pronom « en », d'« appréhender », de « réalité », de « totale » et « immédiate » à la dissertation, puis une première partie qui démontre la thèse au lieu de traiter un seul de ses termes.
\item
  \textbf{Questions transposables.} Connaître ou reconnaître ; l'immédiat et le médiat ; le regard savant contre le regard émerveillé ; la nature comme objet d'usage ou de contemplation : autant de variantes probables du même thème.
\end{itemize}

\end{document}
